\section{The Sifter platform}\label{sec:system}

\subsection{Design goals}

Sifter was built for one recurring situation: a review team has a protocol, a
record set of a few thousand items, and a limited number of reviewer hours.
The platform is the result of four requirements that were treated as
non-negotiable, summarised in Table~\ref{tab:goals}. Each of them rules out a
design that would otherwise have been simpler.

\begin{table}[htbp]
\centering
\caption{Design goals and the mechanism that implements each one in Sifter}
\label{tab:goals}
\begin{tabular}{@{}p{0.19\textwidth}p{0.37\textwidth}p{0.27\textwidth}@{}}
\toprule
Goal & Mechanism & Where it is visible \\
\midrule
Auditability &
One decision row per record, carrying label, justification, model identifier,
prompt version and token counts &
Record detail view; CSV and JSON export \\
Provider independence &
A provider abstraction with one generic OpenAI-compatible adapter plus native
adapters for Gemini, Anthropic and Vertex AI &
Project settings; per-project model choice recorded with the decisions \\
Human primacy &
Three-way decision space in which \emph{doubt} is a first-class label, returned
to the reviewer rather than forced into include or exclude &
Queue filter; doubt tab \\
Measured agreement &
Blind validation mode that hides the model decision until the human decision is
stored, then computes $\kappa$, its confidence interval and the confusion
matrix &
Validation dashboard of the project \\
\bottomrule
\end{tabular}
\end{table}

\subsection{Architecture}

Sifter is a three-tier web application. The backend is a FastAPI service with
SQLAlchemy over PostgreSQL; the frontend is a SvelteKit single-page
application; the whole system is distributed as a Docker Compose stack with
four services (application programming interface, database, test database and
frontend), which is also the installation path recommended to reviewers who
want to run the tool locally. Review data are isolated per project: records,
protocol, decisions and events are scoped to a project identifier, and the
storage layer keeps the associated PDFs in a per-project directory.

The architecture is deliberately unremarkable. The parts that matter for
evidence synthesis are not the web framework but the objects the framework
persists, and those are described next.

\subsection{The protocol object}\label{sec:protocol}

Every project owns a protocol that follows the PICOC structure (population,
intervention, comparison, outcome, context) together with the eligibility
criteria written as inclusion and exclusion rules, the research questions and
the keyword list with synonyms. The protocol is not documentation stored next
to the data: it is the single source from which every screening prompt is
assembled. Changing an eligibility criterion therefore changes the prompt for
every subsequent record, and the version of the protocol used for a given
decision is recorded with that decision.

This decision has a cost and a benefit that should be stated plainly. The cost
is that the platform cannot screen records for which no protocol exists, which
excludes exploratory use on an ad hoc list of papers. The benefit is that the
question ``what exactly was the model asked?'' has a single, inspectable
answer, which is the condition for reproducing the screening at all.

\subsection{Screening pipeline}

Figure~\ref{fig:pipeline} shows the path of a record through the system.
Records enter through RIS or BibTeX import, are deduplicated on DOI and on
normalised title, and are stored with their abstract. For each record the
prompt builder assembles a message from the protocol, the eligibility
criteria and the record text, and the provider adapter sends it to the model
selected for the project. The model is required to return a structured object
containing a label (include, exclude or doubt), a short justification and the
criterion that drove the decision. Responses that do not validate against the
schema are retried, and records that still fail are marked for manual
screening rather than silently dropped.

\begin{figure}[htbp]
\centering
\begin{tikzpicture}[
  stage/.style={draw, rounded corners=2pt, align=center, font=\footnotesize,
                minimum height=10mm, text width=40mm, fill=black!3},
  side/.style={draw, dashed, rounded corners=2pt, align=center, font=\scriptsize,
               minimum height=10mm, text width=34mm, fill=black!1},
  arr/.style={-{Latex[length=2mm]}, semithick}
]
\node[stage] (s1) at (0,0) {Import and deduplication};
\node[stage] (s2) at (0,-1.5) {Prompt assembly};
\node[stage] (s3) at (0,-3.0) {LLM screening};
\node[stage] (s4) at (0,-4.5) {Decision store\\ (include, exclude, doubt)};
\node[stage] (s5) at (0,-6.0) {Human review of doubt and blind sample};
\node[stage] (s6) at (0,-7.5) {Export: PRISMA counts, CSV, LaTeX};

\node[side] (proto) at (5.6,-1.5) {Protocol\\ (PICOC, criteria, prompt version)};
\node[side] (prov)  at (5.6,-3.0) {Provider adapters\\ Gemini, Anthropic,\\ OpenAI-compatible, Vertex AI};
\node[side] (audit) at (5.6,-4.5) {Audit log and cost\\ accounting (tokens per record)};

\draw[arr] (s1) -- (s2);
\draw[arr] (s2) -- (s3);
\draw[arr] (s3) -- (s4);
\draw[arr] (s4) -- (s5);
\draw[arr] (s5) -- (s6);
\draw[arr, dashed] (proto) -- (s2);
\draw[arr, dashed] (prov) -- (s3);
\draw[arr, dashed] (s4) -- (audit);
\end{tikzpicture}
\caption{Path of a record through Sifter. Solid arrows are the screening
pipeline; dashed arrows are the components that make the pipeline
provider-independent and auditable.}
\label{fig:pipeline}
\end{figure}

Two implementation details are worth recording because they change the
interpretation of the results. First, screening is executed in batch: the
platform submits records in parallel and stores one event per state
transition, so the interval between consecutive events measures throughput,
not the time a reviewer would spend reading. Second, the \emph{doubt} label is
not an error state. It is the mechanism through which the model declines to
decide, and the queue gives those records priority over the ones the model
confidently excluded.

\subsection{Blind validation mode}\label{sec:blind}

The validation mode is the feature that distinguishes Sifter from a screening
interface with an API key. When a validation study is opened for a project,
the platform draws a random sample of records, hides the existing model
decision for those records in the reviewer interface, and asks the reviewer to
screen them without seeing the machine label. Only after the human decision is
stored does the interface reveal both decisions and update the statistics:
Cohen's $\kappa$ with a bootstrap confidence interval, raw agreement, the
confusion matrix and the derived sensitivity and specificity, computed against
the human decision as reference \autocite{Cohen_1960,Landis_1977}.

The sample, the seed used to draw it and the resulting pairs are stored, so
the study can be repeated on a different sample and compared with the earlier
one. This is the practical answer to the reproducibility requirement of the
journal's generative AI guidance \autocite{Farotimi_2025}: the agreement study
is an object inside the system, not a table typed into a manuscript.

\subsection{Export and reporting}

The platform exports the record set with both decisions, the audit events and
the counts needed for the PRISMA flow diagram: records identified, duplicates
removed, records screened, records excluded at screening and records sought for
retrieval \autocite{Page_2021}. The export includes the model identifier, the
protocol version and the prompt version for each decision, which is what makes
the screening stage auditable after publication.

\subsection{Implementation and availability}

Version 3.20.0 of the platform is written in Python and TypeScript and is
distributed as a Docker Compose stack. The backend test suite contained 260
tests at the time of the pilot, of which 259 passed and one was skipped, with
75.11\% line coverage against a project gate of 33\%. Source code, licence and
installation instructions are reported in the Data and Code Availability
Statement.
